\chapter{Keeping it running}
\label{ch:operations}

\section{The failure that started this chapter}

The application was deployed on a free plan and worked. Some days later the bot
stopped answering. Nothing had been changed, no error appeared in the logs, and
opening the web page worked normally after a short wait.

The cause was a combination of two behaviours that are documented, free, and
easy to overlook:

\begin{itemize}
  \item the hosting platform \textbf{suspends} the application after about
        fifteen minutes without traffic;
  \item the managed database \textbf{pauses} the project after several days
        without queries.
\end{itemize}

Neither is a bug. Together they break the bot in a way that is invisible: when a
message arrives, Telegram POSTs to a process that is still starting. Telegram
waits a few seconds, gives up, and drops the update. By the time the application
is ready, there is nothing left to answer. The user sees a bot that ignores
them, and the logs show a clean startup.

\begin{keypoint}[The lesson]
A deployment is not finished when the application answers. It is finished when
it still answers a week later without anybody touching it. That last part is a
feature, and it has to be designed like one.
\end{keypoint}

\section{Three pieces}

\begin{table}[h]
\centering
\begin{tabular}{@{}p{3.4cm}p{5.4cm}p{4.7cm}@{}}
\toprule
\textbf{Piece} & \textbf{What it does} & \textbf{Where it runs} \\
\midrule
Uptime monitor & \code{GET /health/db} every five minutes & An external monitor \\
Task cron & \code{POST /tasks/*} on a schedule & An external cron \\
Backup & \code{POST /tasks/backup} weekly & The same cron \\
\bottomrule
\end{tabular}
\caption{What keeps the application alive on a free plan.}
\end{table}

The monitor solves both problems at once. Each request keeps the application
from being suspended, and because the endpoint queries the database, it also
keeps the database from being paused. A single ping every five minutes replaces
both problems with none.

\section{Why the health endpoints differ}

\begin{lstlisting}[style=py]
@app.get("/health")
def health():
    """Liveness: says the process answers. It does not touch the database."""
    return {"status": "ok", "db": "sqlite" if settings.is_sqlite else "postgres"}

@app.get("/health/db")
def health_db():
    """Readiness: a real SELECT 1, and 503 when the database is silent."""
    db = SessionLocal()
    try:
        db.execute(text("SELECT 1"))
    except Exception as exc:
        return JSONResponse({"status": "error", "db": "unreachable",
                             "detail": str(exc)[:200]}, status_code=503)
    finally:
        db.close()
    return {"status": "ok", ...}
\end{lstlisting}

The platform healthcheck points at \code{/health}, which never fails because of
the database. If it did, a five minute database outage would make the platform
consider the application broken and remove the deployment, so a temporary
problem would become a permanent one.

The external monitor points at \code{/health/db}, which fails loudly on purpose.
A paused database with a process still running is exactly the situation that
must produce an alert.

Four tests pin this behaviour, including one that patches the session factory to
raise and asserts that \code{/health} still answers \code{200}.

\section{Scheduled jobs, twice}

The same four jobs can run in two different ways.

On an always-on machine, an in-process scheduler is enough:

\begin{lstlisting}[style=py]
sched.add_job(_refresh_prices, CronTrigger(hour=3, minute=0), id="prices")
sched.add_job(_contributions, CronTrigger(day=1, hour=6, minute=0), id="contributions")
sched.add_job(_snapshot, CronTrigger(hour=23, minute=0), id="snapshot")
if settings.telegram_enabled:
    sched.add_job(_reminder, CronTrigger(day="last", hour=20, minute=0), id="reminder")
\end{lstlisting}

On a host that sleeps it is useless: a scheduler inside a suspended process does
not fire. The same jobs are therefore exposed as HTTP endpoints that an external
cron can call.

\begin{lstlisting}[style=py]
JOBS = {"prices", "contributions", "snapshot", "reminder", "backup"}

@router.post("/{job}")
def run_task(job: str, x_tasks_token: str | None = Header(default=None),
             token: str | None = Query(default=None), db: Session = Depends(get_session)):
    if not _authorised(x_tasks_token, token):
        return JSONResponse({"ok": False, "error": "not authorised"}, status_code=403)
    if job not in JOBS:
        return JSONResponse({"ok": False, "error": "unknown job"}, status_code=404)
    ...
\end{lstlisting}

A single environment variable, \code{USE\_SCHEDULER}, picks the mode. The job
functions are shared, so the two paths cannot drift apart.

\section{Two ways to send the token}

The endpoints accept the token either in a header or in the query string. That
looks like sloppy design, and it is the result of a debugging session.

The header is the better option: it keeps the secret out of URLs, out of browser
history and out of proxy logs. It was the only option in the first version.

Then the cron service refused to work. Every job returned \code{403}, with no
further detail, because a rejected request is all the endpoint is willing to say.
The cause was that the free cron service does not forward custom headers to the
target. The header was simply never arriving.

\begin{lstlisting}[style=py]
def _authorised(header_token, query_token) -> bool:
    expected = settings.tasks_token
    if not expected:               # no token configured: endpoints disabled
        return False
    submitted = header_token or query_token or ""
    return secrets.compare_digest(submitted, expected)
\end{lstlisting}

An empty configured token disables the endpoints entirely, rather than allowing
everything, which is the direction a mistake here should fail in. The comparison
uses \code{compare\_digest} so that a wrong token cannot be discovered by timing
the response.

\begin{pitfall}[Two more traps in the same area]
The cron URL must use \code{https://}. With \code{http://} the platform answers
a redirect, and a redirected POST loses its body, so the job appears to run and
does nothing.

And a job that fails must return a non-2xx status, or the cron dashboard shows
green for a task that never worked. That is why the backup endpoint answers
\code{500} when delivery failed.
\end{pitfall}

\section{Backups}

A free managed database can be paused and, eventually, deleted. The backup
therefore has one rule: the copy must not live inside the thing it protects.

The dump is written by walking the SQLAlchemy metadata, so a new table joins the
backup without anybody editing the backup code:

\begin{lstlisting}[style=py]
EXCLUDED_COLUMNS = {"users": {"password_hash"}}

def export_data(db: Session) -> dict:
    tables = {}
    for name, table in Base.metadata.tables.items():
        excluded = EXCLUDED_COLUMNS.get(name, set())
        columns = [c for c in table.columns if c.name not in excluded]
        rows = db.execute(select(*columns)).all()
        tables[name] = [{c.name: _jsonable(v) for c, v in zip(columns, row)} for row in rows]
    return {"generated": ..., "engine": ..., "tables": tables}
\end{lstlisting}

Three decisions are worth naming.

\textbf{JSON, not \code{pg\_dump}.} The container image is \code{python:3.12-slim}
and carries no Postgres client. Dumping through SQLAlchemy adds no dependency,
produces the same file whether the source is SQLite or Postgres, and can be read
without any tooling.

\textbf{Decimals as strings.} \code{\_jsonable} converts \code{Decimal} to a
string, never to a JSON number, because JSON numbers are floats and the whole
project exists to avoid that.

\textbf{The password hash never leaves.} It is excluded by name, and a test
asserts that no bcrypt prefix appears anywhere in the serialised output. The
user is recreated after a restore with the seed script, so the hash has no value
in the backup and every reason not to travel through a messaging API.

The file is gzipped and sent to the owner's chat, which puts it outside the
provider. Restoring is the mirror image, with one safety rule:

\begin{lstlisting}[style=sh]
python -m scripts.backup restore copy.json.gz            # dry run, prints only
python -m scripts.backup restore copy.json.gz --write    # deletes and inserts
\end{lstlisting}

Restore deletes the contents of every table present in the file, so the
destructive path requires an explicit flag. Tables are deleted in reverse
dependency order and inserted in forward order, using
\code{Base.metadata.sorted\_tables}, so foreign keys hold at every step.
